\begin{statementbox}
	\textbf{\textcolor{CStatement}{条件}}
	\begin{description}[style=nextline, leftmargin=2.8em, labelsep=0.5em, itemsep=0.35em]
		\item[\textbf{\textcolor{CStatement}{条件 1（图形条件）：}}]
			有四个全等的直角三角形，直角边长分别为 $a$ 与 $b$（$a\ge b>0$），斜边长为 $c$。将它们按弦图方式摆放：斜边首尾相接构成边长为 $c$ 的大正方形，中央形成边长为 $a-b$ 的小正方形。
	\end{description}
	\textbf{\textcolor{CStatement}{结论}}
	\begin{description}[style=nextline, leftmargin=2.8em, labelsep=0.5em, itemsep=0.45em]
		\item[\textbf{\textcolor{CStatement}{结论一（勾股定理）：}}]
			\textbf{\textcolor{CStatement}{直观描述：}}
			直角三角形两直角边的平方和等于斜边的平方。
			\[
				a^2 + b^2 = c^2 .
			\]
		\item[\textbf{\textcolor{CStatement}{常见变形（面积恒等式）：}}]
			\textbf{\textcolor{CStatement}{直观描述：}}
			弦图中，斜边上的正方形面积等于四个直角三角形面积与中间小正方形面积之和。
			\[
				c^2 = 4 \times \frac{1}{2}\,ab + (a-b)^2 .
			\]
	\end{description}

\end{statementbox}
